\section{System Model}
\label{sec:system_model}

We consider a downlink \gls{dmimo} system with active \glspl{ap} $\mathcal{A}=\{1,\ldots,M\}$ and scheduled \glspl{ue} $\mathcal{U}=\{1,\ldots,K\}$. Let $N_m$ denote the number of transmit ports at \gls{ap} $m$ and $N_{r,k}$ the number of receive ports at \gls{ue} $k$. The per-\gls{ap} channel block is
\begin{equation}
    \mathbf{H}_{k,m} \in \mathbb{C}^{N_{r,k} \times N_m},
\end{equation}
and the full distributed channel observed by \gls{ue} $k$ is
\begin{equation}
    \mathbf{H}_k = \left[\mathbf{H}_{k,1}\;\mathbf{H}_{k,2}\;\cdots\;\mathbf{H}_{k,M}\right].
\end{equation}

Given a receive combiner $\mathbf{u}_k$, the effective channel becomes
\begin{equation}
    \mathbf{h}_k^{\mathsf{H}} = \mathbf{u}_k^{\mathsf{H}}\mathbf{H}_k.
\end{equation}
Stacking all users yields the effective distributed channel matrix $\mathbf{H}_{\mathrm{eff}}$. The current pipeline applies a least-squares pseudo-inverse \gls{zf} precoder on $\mathbf{H}_{\mathrm{eff}}$, followed by per-stream power allocation under a total \gls{ap} power budget.

For \gls{ue} $k$, the instantaneous downlink \gls{sinr} is evaluated as
\begin{equation}
    \mathrm{SINR}^{\mathrm{DL}}_k =
    \frac{p_k \left|\mathbf{h}_k^{\mathsf{H}}\mathbf{f}_k\right|^2}
    {\sum_{i \in \mathcal{U},\, i \neq k} p_i \left|\mathbf{h}_k^{\mathsf{H}}\mathbf{f}_i\right|^2 + \sigma^2},
\end{equation}
where $\mathbf{f}_i$ is the precoding vector for stream $i$, $p_i$ its allocated power, and $\sigma^2$ the thermal-noise term. The implementation keeps the interference term explicitly and exports desired-signal, interference, and noise power terms for each evaluated user snapshot.

The \gls{esr} used later in the results section is derived from the same per-snapshot \gls{sinr} values. For snapshot $n$ with scheduled-user set $\mathcal{U}_n$, the postprocessing layer computes
\begin{equation}
    \mathrm{ESR}_n = \sum_{k \in \mathcal{U}_n}\log_2\!\left(1+\mathrm{SINR}^{\mathrm{DL}}_{k,n}\right).
\end{equation}
These per-snapshot \gls{sinr} values form the basis of the trajectory-level \gls{sinr} and \gls{esr} statistics reported later in the article.

\subsection{Placement Strategies}

The placement layer compares five strategies under a common \gls{ap} budget:
\begin{itemize}
    \item \texttt{central\_massive\_mimo}: one rooftop proxy \gls{ap} is placed on the building whose rooftop representative point is nearest the area center; the proxy is oriented toward the area center. This is a normalized central reference, not a full single co-located-array placement model.
    \item \texttt{distributed\_fixed}: a distributed set of wall-mounted candidate positions is sampled once and kept fixed over the full trajectory.
    \item \texttt{distributed\_movable}: optimization 1. The same number of distributed \glspl{ap} is reselected at each relocation window from the wall-mounted candidate pool with a history-aware local-\gls{csi} rule.
    \item \texttt{distributed\_movable\_optimization\_2}: optimization 2. The same number of distributed \glspl{ap} is reselected at each relocation window from the wall-mounted candidate pool with a current-window nearest-user \gls{ue}-to-\gls{ue} proxy-\gls{csi} rule and a proxy-\gls{esr} objective.
    \item \texttt{distributed\_movable\_optimization\_3}: optimization 3. The same number of distributed \glspl{ap} is reselected at each relocation window from the wall-mounted candidate pool with the same nearest-user proxy association as optimization 2, but with a proxy-average-power objective based on $|\mathrm{CSI}|^2$.
\end{itemize}

Let $\mathcal{C}$ denote the wall-mounted candidate pool and let $\mathcal{S}_w \subseteq \mathcal{C}$ denote the selected movable-\gls{ap} subset in relocation window $w$. In the Rabot scenario studied here, $\lvert \mathcal{S}_w \rvert = N_{\mathrm{mov}}=6$ and no always-on distributed \glspl{ap} are used, so the distributed comparison isolates the effect of changing the movable subset over time.

\subsection{Relocation Heuristic}

At an intuitive level, the relocation mechanism repeats the same observe-score-relocate loop at every relocation instant. Optimization~1 uses simulated \gls{ap}-to-\gls{ue} channel observations from previously evaluated subsets and discounts older samples exponentially. Optimization~2 and optimization~3 follow a different practical assumption: the \gls{ap}-to-\gls{ue} channel of a dormant candidate position is not known before an \gls{ap} is actually moved there. Consequently, these two rules do not optimize directly over candidate-position \gls{ap}-to-\gls{ue} \gls{csi}. Instead, each candidate site is associated with the closest user within a distance threshold, the measured \gls{ue}-to-\gls{ue} channel of that nearby user is reused as a proxy for the dormant candidate site, and the subset is ranked either by proxy \gls{esr} (optimization~2) or by proxy average received power (optimization~3). Once the best subset has been chosen, the movable \glspl{ap} are relocated to those candidate positions and the actual deployed \gls{ap}-to-\gls{ue} \gls{csi} is recomputed for storage, postprocessing, and final reporting.

The movable distributed strategy is implemented as a windowed receding-horizon controller. Let $\{t_n\}_{n=0}^{N_t-1}$ denote the simulated snapshot times and let $T_{\mathrm{rel}}$ denote the relocation interval. The trajectory is partitioned into disjoint decision windows
\begin{equation}
    \mathcal{W}_w = \left\{n \in \{0,\ldots,N_t-1\} :
    \left\lfloor \frac{t_n-t_0}{T_{\mathrm{rel}}}\right\rfloor = w\right\}.
\end{equation}
At the beginning of each window $w$, the controller selects a subset $\mathcal{S}_w \subseteq \mathcal{C}$ with cardinality $\lvert \mathcal{S}_w \rvert = N_{\mathrm{mov}}$, where $N_{\mathrm{mov}}$ is the movable-\gls{ap} budget. The fixed distributed mode uses one time-invariant subset, whereas the movable mode recomputes $\mathcal{S}_w$ at every relocation instant.

\paragraph{Optimization 1: historical $K$-nearest local-\gls{csi} controller.}
The first online decision rule is intentionally local and does not optimize a full radio map. For a candidate site $c \in \mathcal{C}$ with horizontal position $\mathbf{x}_c \in \mathbb{R}^2$, define the local neighborhood $\mathcal{N}_K(c,w)$ as the set of the $K$ nearest \gls{ue}-snapshot pairs in the horizontal plane within window $w$,
\begin{equation}
    \mathcal{N}_K(c,w) =
    \operatorname{KNN}_K\!\left(
        \mathbf{x}_c;
        \{\mathbf{r}_{u,n} : u \in \mathcal{U}, n \in \mathcal{W}_w\}
    \right),
\end{equation}
where $\mathbf{r}_{u,n}$ denotes the horizontal position of \gls{ue} $u$ at time index $n$. For a subset $\mathcal{S} \subseteq \mathcal{C}$, the support region used by the controller is the union
\begin{equation}
    \mathcal{N}_K(\mathcal{S},w) = \bigcup_{c \in \mathcal{S}} \mathcal{N}_K(c,w).
\end{equation}
In the Rabot configuration, $K=10$, so each candidate contributes at most ten local \gls{ue}-snapshot pairs per window. This restriction makes the relocation rule sensitive to the \gls{csi} observed near the candidate sites rather than to globally averaged scene statistics. In other words, user positions are used to construct the local sampling mask, but they do not enter the score directly.

For every queried subset $\mathcal{S}$, the implementation computes the corresponding \gls{ap}-to-\gls{ue} \gls{csi}/\gls{sinr} segment over the current window and stores it in a per-subset cache. Thus, the relocation heuristic is driven by channel-response information associated with the candidate subset under evaluation, not by geometry alone. If $\gamma_{u,n}(\mathcal{S})$ denotes the best-link \gls{sinr} sample produced when subset $\mathcal{S}$ is active, the historical record available when deciding window $w$ is
\begin{equation}
    \mathcal{H}_w(\mathcal{S}) =
    \left\{(t_n,\gamma_{u,n}(\mathcal{S})) :
    u \in \mathcal{U}, n \in \mathcal{W}_j, j = 0,\ldots,w\right\}.
\end{equation}
When the same subset is revisited later, previously computed window segments are reused and only the new current-window segment is appended. Thus, the historical objective remains consistent while avoiding repeated ray-tracing evaluations for past windows. In the present implementation, the relocation score is formed from these \gls{ap}-to-\gls{ue} channel samples only; \gls{ue}-to-\gls{ue} \gls{csi} is reserved for the final trajectory-level comparison score, and \gls{ap}-to-\gls{ap} \gls{csi} is exported for analysis rather than for online relocation control.

Historical CSI is discounted through exponential aging. If a sample was observed at time $t_n$ and the relocation decision for window $w$ is taken at $t_w^{\mathrm{end}}$, the associated weight is
\begin{equation}
    \omega_n^{(w)} = \exp\!\left[-\lambda\left(t_w^{\mathrm{end}}-t_n\right)\right],
\end{equation}
with $\lambda$ the configured decay rate. Hence, newer samples always dominate older ones: a sample measured at the decision time keeps weight $1$, a sample that is one relocation window old keeps weight $\exp(-\lambda T_{\mathrm{rel}})$, and the weight keeps halving every additional relocation window when $\lambda=\ln(2)/T_{\mathrm{rel}}$, as in the Rabot configuration. The local control objective is then the weighted empirical 10th percentile of the \gls{sinr} samples retained by the local neighborhoods across all observed windows,
\begin{equation}
    J_w^{(1)}(\mathcal{S}) =
    Q_{0.1}^{\omega}\!\left(
        \left\{\gamma_{u,n}(\mathcal{S}) :
        (u,n) \in \bigcup_{j=0}^{w} \mathcal{N}_K(\mathcal{S},j)\right\}
    \right),
\end{equation}
where $Q_{0.1}^{\omega}(\cdot)$ denotes the weighted 10th-percentile operator induced by $\omega_n^{(w)}$. The use of a lower-tail percentile, rather than a mean, makes the controller deliberately conservative with respect to persistently weak local links.

\paragraph{Optimization 2: current-window nearest-user proxy-\gls{csi} controller.}
The second online rule removes the $K$-nearest construction and does not use historical samples. Let $d_{\mathrm{th}}$ denote the configured distance threshold. For each candidate site $c \in \mathcal{C}$ and snapshot $n \in \mathcal{W}_w$, define the anchor user
\begin{equation}
    \ell(c,n)=
    \arg\min_{u \in \mathcal{U}}
    \left\|\mathbf{r}_{u,n}-\mathbf{x}_c\right\|_2,
\end{equation}
and the validity indicator
\begin{equation}
    \delta(c,n)=
    \mathbf{1}\!\left\{
        \left\|\mathbf{r}_{\ell(c,n),n}-\mathbf{x}_c\right\|_2 \le d_{\mathrm{th}}
    \right\}.
\end{equation}
The dormant candidate-position \gls{ap}-to-\gls{ue} channel is not available before relocation. Therefore, optimization~2 replaces it by a proxy channel extracted from the measured \gls{ue}-to-\gls{ue} \gls{csi}. If $h^{\mathrm{UU}}_{u,v,n}$ denotes the narrowband \gls{ue}-to-\gls{ue} coefficient from transmitter-user $v$ to receiver-user $u$ at snapshot $n$, the proxy coefficient used for candidate $c$ is
\begin{equation}
    \tilde{h}_{u,c,n}=
    \begin{cases}
        h^{\mathrm{UU}}_{u,\ell(c,n),n}, & \delta(c,n)=1,\; u \neq \ell(c,n), \\
        h^{\mathrm{UU}}_{v^\star,\ell(c,n),n}, & \delta(c,n)=1,\; u = \ell(c,n), \\
        0, & \delta(c,n)=0,
    \end{cases}
\end{equation}
where $v^\star=\arg\max_{v \in \mathcal{U}} |h^{\mathrm{UU}}_{v,\ell(c,n),n}|$ supplies a surrogate self-link because the diagonal user-to-itself measurement is not observable. For a candidate subset $\mathcal{S}$, the proxy channel matrix $\tilde{\mathbf{H}}_n(\mathcal{S})$ is formed by stacking the columns $\tilde{h}_{u,c,n}$ for all $c \in \mathcal{S}$. The same \gls{zf} procedure as in the deployed \gls{ap} case is then applied to $\tilde{\mathbf{H}}_n(\mathcal{S})$, yielding a proxy \gls{sinr} $\tilde{\gamma}_{u,n}(\mathcal{S})$ and the current-window proxy-\gls{esr} objective
\begin{equation}
    J_w^{(2)}(\mathcal{S}) =
    \frac{1}{|\mathcal{W}_w|}
    \sum_{n \in \mathcal{W}_w}
    \sum_{u \in \mathcal{U}}
    \log_2\!\left(1+\tilde{\gamma}_{u,n}(\mathcal{S})\right).
\end{equation}
Consequently, optimization~2 uses measured \gls{ue}-to-\gls{ue} proxy \gls{csi} to rank candidate subsets, while avoiding the unrealistic assumption that dormant candidate-position \gls{ap}-to-\gls{ue} channels are already known.

\paragraph{Optimization 3: current-window nearest-user proxy-power controller.}
The third online rule reuses the same nearest-user anchor $\ell(c,n)$ and validity indicator $\delta(c,n)$ as optimization~2, but it does not use the complex proxy coefficient directly. Instead, it works with the proxy received-power quantity
\begin{equation}
    \tilde{\rho}_{u,c,n} =
    \frac{P_{\mathrm{AP}}}{P_{\mathrm{UE}}}
    \left|\tilde{h}_{u,c,n}\right|^2,
\end{equation}
where the factor $P_{\mathrm{AP}}/P_{\mathrm{UE}}$ rescales the measured \gls{ue}-to-\gls{ue} proxy power to the distributed-\gls{ap} transmit-power budget. The optimization-3 objective is then
\begin{equation}
    J_w^{(3)}(\mathcal{S}) =
    \frac{1}{|\mathcal{W}_w|\,|\mathcal{U}|}
    \sum_{n \in \mathcal{W}_w}
    \sum_{u \in \mathcal{U}}
    \sum_{c \in \mathcal{S}}
    \tilde{\rho}_{u,c,n}.
\end{equation}
Optimization~3 is therefore the power-only counterpart of optimization~2: it uses the same nearest-user proxy association, but it requires only $|\mathrm{CSI}|^2$ rather than the complex proxy channel.

\paragraph{Greedy subset construction.}
All three optimization rules use the same forward-selection search. Starting from $\mathcal{S}_w^{(0)}=\varnothing$, the controller adds one candidate at a time according to
\begin{equation}
    c_r^\star =
    \arg\max_{c \in \mathcal{C}\setminus \mathcal{S}_w^{(r-1)}}
    J_w^{(q)}\!\left(\mathcal{S}_w^{(r-1)} \cup \{c\}\right),
\end{equation}
\begin{equation}
    q \in \{1,2,3\},
    \qquad
    \mathcal{S}_w^{(r)} =
    \mathcal{S}_w^{(r-1)} \cup \{c_r^\star\},
    \qquad r=1,\ldots,N_{\mathrm{mov}},
\end{equation}
and finally sets $\mathcal{S}_w=\mathcal{S}_w^{(N_{\mathrm{mov}})}$. Consequently, the current implementation is a forward-selection heuristic rather than an exhaustive combinatorial optimizer.

After either optimization rule has chosen the subset, the selected geometric candidate sites are converted into physical relocations of labeled movable \glspl{ap}. Let $\mathbf{p}_m^{(w-1)}$ denote the previous-window position of movable \gls{ap} $m$ and let $\mathbf{q}_j^{(w)}$ denote the coordinates of the candidates in the newly selected subset. The controller solves the linear assignment problem
\begin{equation}
    \pi_w^\star =
    \arg\min_{\pi \in \Pi_{N_{\mathrm{mov}}}}
    \sum_{m=1}^{N_{\mathrm{mov}}}
    \left\|\mathbf{p}_m^{(w-1)}-\mathbf{q}_{\pi(m)}^{(w)}\right\|_2^2,
\end{equation}
where $\Pi_{N_{\mathrm{mov}}}$ is the permutation set on $\{1,\ldots,N_{\mathrm{mov}}\}$. This minimum-displacement matching preserves \gls{ap} identities across windows and yields schedule traces that correspond to physically meaningful AP trajectories instead of arbitrary candidate relabelings.

Once all windows have been simulated, the concatenated movable-strategy trajectory is evaluated using the same trajectory-level summary score as the static baselines,
\begin{equation}
    F = -\alpha P_{\mathrm{out}} + \beta P_{10} + \eta T_{\mathrm{peer}},
\end{equation}
where $P_{\mathrm{out}}$ is the fraction of user snapshots below the configured \gls{sinr} threshold, $P_{10}$ is the trajectory-wide 10th percentile, and $T_{\mathrm{peer}}$ is a peer-weighted mean \gls{sinr} term derived from \gls{ue}-to-\gls{ue} \gls{csi}. The coefficients $(\alpha,\beta,\eta)$ correspond to the configured outage, percentile, and peer tie-break weights. This final score is used for reporting and mode comparison; the online relocation controller itself is driven by one of $J_w^{(1)}(\mathcal{S})$, $J_w^{(2)}(\mathcal{S})$, or $J_w^{(3)}(\mathcal{S})$, depending on the selected movable-strategy variant.
The above formulation defines the relocation mechanisms used throughout the Rabot simulations.
